\Needspace{6\baselineskip}
\subsection{Reference Appliance and System Components}

The system was developed for the Rheem EW1-50RP electric storage water heater, which has a nominal capacity of 50 L
and an electrical power rating of 1500 W at 220--240 V. The appliance includes an internal temperature controller,
a digital display, and physical buttons for operating the heating system~\cite{rheemmanual}.

The proposed IoT architecture extends the appliance's monitoring and control capabilities through external sensing,
mechanical actuation, and wireless communication. An ESP32-WROOM microcontroller serves as the central embedded
controller, while an ESP32-CAM operates as an independent image acquisition node.

Table~\ref{tab:components} summarizes the components and their corresponding roles within the system.

\begin{table*}[htbp]
\centering
\caption{Main hardware and software components of the IoT system.}
\label{tab:components}
\small
\renewcommand{\arraystretch}{1.12}
\begin{tabularx}{\textwidth}{@{}P{2.7cm} P{2.8cm} L@{}}
\toprule
\tblhead{Function} & \tblhead{Component} & \tblhead{Description} \\
\midrule

Central controller
& ESP32-WROOM-32
& Coordinates sensor acquisition, servo actuation, and MQTT communication over Wi-Fi. \\

Display capture
& ESP32-CAM / OV2640
& Serves JPEG images of the original display over HTTP. \\

Pilot sensing
& TCS34725 RGB sensor
& Acquires raw RGBC readings from the pilot through I\textsuperscript{2}C. \\

Flow acquisition
& SEA ZJ-S201 Hall-effect flow sensor
& Provides the pulse-count input of the separate flow channel; a configured factor converts counts into flow rate and accumulated volume. \\

Mechanical actuation
& Low-voltage servomotors
& Press and release the original panel buttons. \\

Wireless communication
& Wi-Fi and MQTT
& Enable the exchange of control commands and telemetry between the ESP32 and application infrastructure. \\

Image processing
& Sharp, seven-segment reader, Tesseract.js
& Preprocess images and provide configurable numerical readers for the original display. \\

Application service
& Node.js and WebSocket
& Manage image processing, MQTT messaging, application state, and bidirectional communication with the web interface. \\

User interface
& Web-based dashboard
& Presents device information and sensor measurements and provides controls for transmitting commands to the system. \\

\bottomrule
\end{tabularx}
\end{table*}

\Needspace{6\baselineskip}
\subsection{System Architecture}

The system follows a distributed architecture composed of an embedded control subsystem, an image acquisition subsystem,
a communication layer, and a web-based application.
Figure~\ref{fig:architecture} illustrates the principal components and the communication paths between them.

\begin{figure*}[!t]
    \centering
    \resizebox{0.96\linewidth}{!}{\resizebox{\textwidth}{!}{%
\begin{tikzpicture}[
  font=\small,
  box/.style={draw, rounded corners, align=center, text width=3.4cm, minimum height=1.15cm, inner sep=6pt},
  link/.style={-{Latex[length=2mm]}, thick},
  labeltext/.style={font=\footnotesize, align=center, fill=white, inner sep=2pt}
]
\node[box, text width=4cm] (heater) at (0,0) {Water heater panel\\Numeric display and pilot\\Power / increase / decrease};
\node[box] (camera) at (-4,-2.5) {ESP32-CAM\\Display image capture};
\node[box] (rgb) at (0,-2.5) {RGB sensor\\Pilot light acquisition};
\node[box] (servos) at (4,-2.5) {Three servos\\Power / increase / decrease};
\node[box] (backend) at (-4,-5.4) {Application backend\\Image preprocessing / OCR\\MQTT client};
\node[box, text width=1.15cm, inner sep=4pt] (broker) at (-1,-5.4) {MQTT\\broker};
\node[box] (controller) at (2,-5.4) {ESP32 controller\\RGB input / flow pulse ISR\\Queued actuator control};
\node[box] (dashboard) at (-4,-8) {Dashboard\\Measurements and requests};
\node[box, dashed] (flow) at (2,-8) {Separate flow channel\\Hall-effect pulse sensor};

\draw[link] (heater.west) -| node[labeltext, pos=0.75, left] {Display\\image} (camera.north);
\draw[link] (heater.south) -- node[labeltext, right] {Pilot light} (rgb.north);
\draw[link] (servos.north) |- node[labeltext, pos=0.3, right] {Mechanical\\button press} (heater.east);
\draw[link] (camera.south) -- node[labeltext, right] {Wi-Fi / HTTP\\JPEG capture} (backend.north);
\draw[link] (rgb.south) -- node[labeltext, right] {I\textsuperscript{2}C} (controller.north);
\draw[link] (controller.east) -| node[labeltext, pos=0.75, right] {PWM} (servos.south);
\node[labeltext] at (-1,-4.3) {Wi-Fi / MQTT\\Telemetry and commands};
\draw[thick, {Latex[length=2mm]}-{Latex[length=2mm]}] (backend.east) -- (broker.west);
\draw[thick, {Latex[length=2mm]}-{Latex[length=2mm]}] (broker.east) -- (controller.west);
\draw[thick, {Latex[length=2mm]}-{Latex[length=2mm]}] (backend.south) -- node[labeltext, right] {WebSocket} (dashboard.north);
\draw[link, dashed] (flow.north) -- node[labeltext, right] {Pulse input / ISR} (controller.south);
\end{tikzpicture}%
}
}
    \caption{Prototype architecture. The camera-to-server recognition path adapts the optical acquisition principle of Tongbai~\cite{tongbai2024}; the control, pilot, and separate flow paths are the implementation developed here.}
    \label{fig:architecture}
\end{figure*}

The ESP32-WROOM operates as the central embedded controller and interfaces with the sensors and servomotors.
The TCS34725 communicates through the I\textsuperscript{2}C protocol, while the flow sensor provides digital
pulses that are counted through a GPIO interrupt.
These measurements are processed periodically to obtain sensor readings and water consumption metrics.

The ESP32-CAM operates as an independent Wi-Fi node.
Instead of transmitting camera data directly to the ESP32 controller, it exposes an HTTP endpoint through which
the application service retrieves images of the heater's display.
This separation allows image processing to be performed outside the embedded control loop.

The communication layer uses an MQTT broker to exchange messages between the ESP32 and the application service.
The embedded controller subscribes to command topics and publishes telemetry and device state messages.
The application service receives these messages and communicates with the web interface through WebSocket.

User interactions generate WebSocket messages that are processed by the application service and translated into MQTT commands.
These commands are received by the ESP32, which executes the corresponding control routines.

\Needspace{6\baselineskip}
\subsection{Temperature Acquisition and Processing}

Temperature information is acquired from the heater's original digital display using the ESP32-CAM.
The camera provides JPEG images through its HTTP interface, which are requested and processed by the Node.js application service.

The image processing pipeline uses Sharp to perform operations such as cropping, resizing, grayscale conversion, and intensity adjustments.
The selected operations align and isolate the digits before recognition.

The application selects either a geometric seven-segment reader or Tesseract.js to identify the numerical characters. The geometric reader compares the illuminated segment pattern with digit patterns; the Tesseract route applies its configured text-confidence threshold.
The recognized value is validated before being transmitted through MQTT using the topic \texttt{heater/input/temperature}.

The ESP32 receives the temperature value, updates its internal representation, and publishes the corresponding state through
\texttt{heater/state/temperature}. This information is then transmitted to the web interface.

\Needspace{6\baselineskip}
\subsection{Water Flow and Indicator Monitoring}

The separate flow-acquisition channel connects a Hall-effect turbine sensor to the ESP32.
Each pulse generated by the sensor is registered through an interrupt-based counting mechanism.

For a measurement interval, the water volume is obtained from the pulse count and the configured pulses-per-liter calibration factor:

\begin{equation}
    V = \frac{N}{K},
    \label{eq:water-volume}
\end{equation}

where $V$ represents the water volume in liters, $N$ is the number of registered pulses, and $K$ is the calibration factor
expressed in pulses per liter.

The average flow rate over the measurement interval is calculated as:

\begin{equation}
    Q = \frac{60V}{\Delta t},
    \label{eq:flow-rate}
\end{equation}

where $Q$ is expressed in liters per minute and $\Delta t$ represents the measurement interval in seconds.
The accumulated volume is updated by adding the volume obtained during each interval.

The TCS34725 sensor acquires RGB and clear-channel measurements from the heater's indicator.
Communication with the ESP32 is established through I\textsuperscript{2}C.
The acquired values provide optical information about the indicator and are included in the system's telemetry.

\Needspace{6\baselineskip}
\subsection{Communication and Application Integration}
A Mosquitto broker distributes MQTT commands and telemetry. The ESP32 subscribes to validated control requests, including \texttt{heater/set/power} and \texttt{heater/set/autoshutdown}, and publishes readings through state topics. The Node.js service combines this MQTT client with HTTP image acquisition and a WebSocket server. The dashboard uses that WebSocket connection for requests and status updates, keeping image recognition separate from the embedded control loop.
